% ============================================================
%  I. INTRODUCTION
% ============================================================
\section{Introduction}
\label{sec:introduction}

\IEEEPARstart{F}{ocused} \acf{FAST} is a standardised bedside ultrasound
protocol that enables clinicians to rapidly detect life-threatening internal
injuries---haemopericardium, haemothorax, and free intraperitoneal
fluid---at the point of care~\cite{rozycki1996}. Its speed and
non-invasiveness make it indispensable in trauma bays, pre-hospital emergency
medicine, and military field operations. However, \ac{FAST} requires a trained
sonographer who is physically present and manually positions the probe at four
canonical anatomical windows: the subxiphoid, right upper quadrant, left upper
quadrant, and suprapubic regions. In mass-casualty incidents, pandemics, or
remote deployment scenarios, qualified personnel are often unavailable or
overwhelmed, creating a critical bottleneck in triage.

Autonomous robotic ultrasound has the potential to address this gap. A
comprehensive review of robot-assisted ultrasound~\cite{du2024review} and
recent fixed-base systems~\cite{mathiassen2016, li2021overview, duan2022}
demonstrate that high-quality probe positioning can be achieved without direct
human guidance. In our own prior work~\cite{farsoni2025autonomous}, we showed
that a collaborative robot arm on a fixed base can autonomously screen the
abdominal aorta using impedance control and deep learning segmentation,
achieving clinically relevant diameter measurements. However, all these systems
share a fundamental limitation: they require the patient to be pre-positioned
within the arm's reachable workspace, confining autonomous ultrasound to
controlled clinical settings.

Legged platforms such as the Boston Dynamics Spot offer a qualitative step
beyond this limitation. They can navigate stairs, rubble, and confined spaces
inaccessible to wheeled robots, and can therefore reach patients wherever they
lie. Mounting a serial manipulator arm on a legged base introduces, however, a
fundamental control challenge: the arm workspace is small relative to the
environment, so the base must actively cooperate with the arm throughout the
mission. Simply stopping the base while the arm operates---the prevalent
approach in mobile manipulation~\cite{bellicoso2019}---wastes time and fails
when the target falls outside the reachable workspace. \ac{WBC} formulations
that simultaneously optimise base and arm motion~\cite{sleiman2021,
haviland2022} address this, but have not yet been applied to
perception-driven medical scanning tasks where the patient location is unknown
and must be estimated on-the-fly from noisy visual detections.

This paper presents \acf{TERESA}, a complete system for autonomous \ac{FAST}
scanning on a Boston Dynamics Spot legged robot equipped with a Unitree Z1
manipulator arm. Building on the fixed-base scanning competence established
in~\cite{farsoni2025autonomous}, we extend autonomous ultrasound to a fully
mobile platform capable of operating in unstructured environments. The core of
the system is a whole-body controller that formulates base and arm motion as a
single \ac{QP} over a holistic $6{\times}8$ Jacobian, avoiding the
stop-then-manipulate paradigm. A dual-camera perception pipeline---wide-field
Orbbec \ac{RGBD} for navigation and \ac{EE}-mounted Intel RealSense for
surface estimation---feeds a six-state \ac{FSM} that orchestrates the full
mission lifecycle from patient detection to completed scan.

The main contributions of this paper are:
\begin{enumerate}
  \item A \textbf{\ac{WBC} formulation} for a non-holonomic legged base
        coupled with a six-\ac{DOF} arm, with manipulability-weighted effort
        redistribution between base and arm (Section~\ref{sec:method}).
  \item A \textbf{dual-camera perception pipeline} that provides robust
        patient localisation across the full approach distance, with a
        principled sensor handoff from wide-field skeleton tracking to
        \ac{EE} surface estimation (Section~\ref{sec:system_architecture}).
  \item A \textbf{six-state mission \ac{FSM}} that integrates \ac{WBC},
        perception, impedance contact control, and reactive workspace
        extension into a unified autonomous scanning system
        (Section~\ref{sec:system_architecture}).
  \item \textbf{Experimental validation} on a physical Spot+Z1 platform
        across a $3{\times}3$ grid of patient positions and orientations,
        demonstrating a \XX\,\% reduction in mission time and probe
        positioning error below \XX\,mm compared to a stop-and-manipulate
        baseline (Section~\ref{sec:results}).
\end{enumerate}

To the best of our knowledge, this is the first system to demonstrate fully
autonomous \ac{FAST} scanning on a legged mobile manipulator operating in
unstructured environments.

The remainder of this paper is organised as follows.
Section~\ref{sec:related_work} reviews related work in robotic ultrasound,
mobile manipulation, and legged robot control.
Section~\ref{sec:problem_formulation} formally defines the system and the
problem. Section~\ref{sec:method} derives the \ac{WBC} formulation.
Section~\ref{sec:system_architecture} describes the complete system
architecture. Section~\ref{sec:experiments} details the experimental setup
and Section~\ref{sec:results} presents the results and discussion.
Section~\ref{sec:conclusion} concludes the paper.
